\section{课程要求与代理可交付材料}
\small
\begin{longtable}{p{0.22\textwidth}p{0.43\textwidth}p{0.27\textwidth}}
\toprule 课程要求&本次可核查技术材料&待人工完成或适用限制\\\midrule
\endfirsthead\toprule 课程要求&本次可核查技术材料&待人工完成或适用限制\\\midrule\endhead
DRAM、SEE与AI方法&第1至4章、公式及参数来源&学生个人理解与答辩表现\\
2至3人合作&可按模型、平台、数据分析登记职责&姓名、学号、真实贡献和日期\\
可行性、方案、步骤和进度&第1、4、5章与materials索引&教师对方案和变更的审核\\
环境与资源&Windows和Linux VM的真实远程访问式流程&实验室服务器/模型正式交接\\
基础模型与源项传递&二维1T1C输入、G4事件、源积分审计&未校准真实DRAM工艺\\
自然语言参数遍历&R50已验收preview/audit/locate及拒绝行为&学生现场演示尚未发生\\
多尺寸/粒子/方向评估&完整64事件、128响应与固定分层&范围外的能量、位置与工艺未验证\\
拟合、外推与验证&训练LOCO、64冻结预测及全量独立评分&未建立商业可靠性外推\\
工艺与数值敏感性&14氧化层case、固定两边界目标及匹配细化&未完成变化材料栈的输运重算\\
总结与作品演示&完整报告、研究稿、PPT及只读演示脚本&教师评审、试讲、答辩、考核待发生\\
\bottomrule\end{longtable}\normalsize
\section{成员与实际活动登记}
指导教师李洋来自课程大纲，不能据此声称本报告已获教师审核。
真实参与者应填写成员姓名/学号、职责、完成时间与可核查贡献，并另存方案审核、
资源交接、试讲反馈、正式演示/答辩及教师考核记录。AI工具不计作小组成员。
课程评分权重为导师20\%、报告20\%、答辩或作品演示60\%；本报告不预测成绩。

\section{只读复现与演示}
\begin{enumerate}
\item 核验本包\sourcepath{sources.json}、\sourcepath{delivery-manifest.json}及输入快照清单，说明固定分母和模型边界。
\item 展示受约束R50请求的预览、审计与指定事件定位。历史映射快照的24小时语法字段不决定当前无限预算。
\item 选择一行\sourcepath{data/snapshot/response-sources.json}，定位actual tag、owner plan、case-freeze、run ID和input/archive SHA。
\item 对照源事件和匹配control/neutral，解释2.8~ns的$\Delta Q$及4~ns的$D,M,\Delta M$。
\item 展示预测冻结时间及真实测试提交索引，分别解释留出/外推误差、原r4评分和最终数值证据评分。
\item 展示DATA PLOT、原生场图和owned窗口的不同证据用途；核场变量、单位、4~ns和run来源。
\item 解释R50主线四个原始FAILED、null退出码及相应确定性重试；完整历史账本7条FAILED均保留，费用以完整唯一账本为准。
\end{enumerate}
上述是供参与者采用的步骤，不是已完成学生现场演示记录。耗时求解使用事先冻结归档并明确标注预先运行；
本报告构建及演示资料均不启动新仿真。
\par\medskip\noindent
\begin{tabular}{lp{0.62\textwidth}}
\toprule
证据类别 & 本包入口\\
\midrule
完整128槽及数值状态 & \sourcepath{data/snapshot/responses.csv}\\
真实case owner与四alias & \sourcepath{data/snapshot/response-sources.json}\\
原FAILED和费用 & \sourcepath{data/snapshot/ledger-snapshot.json}\\
原r4固定模型评分 & \sourcepath{data/snapshot/original-r4-prediction-scores.json}\\
最终数值证据评分 & \sourcepath{data/snapshot/final-numerical-prediction-scores.json}\\
原生PNG真实目视 & \sourcepath{data/native-visual-review.json}\\
全部来源与哈希 & \sourcepath{sources.json}\\
\bottomrule
\end{tabular}

